\chapter{Verificación de correo electrónico con OAuth2}
\label{ch:verificacion_correo}

\section{Objetivo}

Implementar un flujo de verificación de correo electrónico para el registro de nuevos usuarios, garantizando que:

\begin{itemize}[leftmargin=*]
  \item Solo se registren correos reales y activos.
  \item El envío de correos funcione de forma autónoma, sin depender de la habilitación de contraseñas de aplicación por parte del administrador de TI de la universidad.
  \item El sistema utilice OAuth2 en lugar de contraseñas SMTP tradicionales, asegurando su viabilidad a largo plazo.
  \item Se mantenga un atajo seguro para crear usuarios de prueba desde la terminal durante el desarrollo.
\end{itemize}

\section{Contexto}

El flujo de registro inicial solo validaba el dominio \texttt{@cua.uam.mx} de forma sintáctica, lo que permitía registrar correos inexistentes o mal escritos. Para enviar correos de verificación reales, se intentó usar contraseñas de aplicación de Gmail, pero el administrador de Google Workspace de la UAM-C tiene esa opción deshabilitada por políticas de seguridad.

Se optó por implementar OAuth2, que permite al backend autenticarse con los servidores de Google usando un \textit{Refresh Token} de larga duración, sin necesidad de contraseñas de aplicación y sin intervención del administrador.

\section{Desarrollo}

\subsection{Configuración en Google Cloud Console}

Se creó un proyecto en Google Cloud Console y se habilitó la Gmail API. Se generó un \textbf{ID de cliente de OAuth} de tipo \textit{Aplicación web} (necesario para el OAuth 2.0 Playground) y se añadió la URI de redirección:

\begin{lstlisting}[language=bash]
https://developers.google.com/oauthplayground
\end{lstlisting}

Utilizando el Playground, se autorizó el scope \texttt{https://mail.google.com/} y se obtuvo un \textbf{Refresh Token} de larga duración.

\subsection{Configuración de Spring Boot}

Se añadieron las dependencias necesarias al \texttt{pom.xml}:

\begin{lstlisting}[language=XML]
<dependency>
    <groupId>org.springframework.boot</groupId>
    <artifactId>spring-boot-starter-mail</artifactId>
</dependency>
<dependency>
    <groupId>com.google.api-client</groupId>
    <artifactId>google-api-client</artifactId>
    <version>2.7.0</version>
</dependency>
\end{lstlisting}

Las credenciales se almacenaron en el archivo \texttt{.env} (ignorado por Git):

\begin{lstlisting}[language=bash]
GMAIL_CLIENT_ID=tu-client-id-web.apps.googleusercontent.com
GMAIL_CLIENT_SECRET=tu-client-secret-web
GMAIL_REFRESH_TOKEN=tu-refresh-token
GMAIL_FROM_EMAIL=tu-correo@gmail.com
\end{lstlisting}

\subsection{Componentes del backend}

\subsubsection{\texttt{MailConfig}}

Clase de configuración que construye el \texttt{JavaMailSender} autenticándose mediante XOAUTH2. Renueva el token de acceso a partir del Refresh Token antes de cada envío.

\begin{lstlisting}[language=Java]
@Configuration
public class MailConfig {
    @Value("${GMAIL_CLIENT_ID}") private String clientId;
    @Value("${GMAIL_CLIENT_SECRET}") private String clientSecret;
    @Value("${GMAIL_REFRESH_TOKEN}") private String refreshToken;
    @Value("${spring.mail.username}") private String username;

    @Bean
    public JavaMailSender javaMailSender() throws Exception {
        JavaMailSenderImpl mailSender = new JavaMailSenderImpl();
        mailSender.setHost("smtp.gmail.com");
        mailSender.setPort(587);
        mailSender.setUsername(username);

        GoogleCredential credential = new GoogleCredential.Builder()
                .setTransport(GoogleNetHttpTransport.newTrustedTransport())
                .setJsonFactory(GsonFactory.getDefaultInstance())
                .setClientSecrets(clientId, clientSecret)
                .build()
                .setRefreshToken(refreshToken);

        credential.refreshToken();
        mailSender.setPassword(credential.getAccessToken());

        Properties props = mailSender.getJavaMailProperties();
        props.put("mail.smtp.auth", "true");
        props.put("mail.smtp.starttls.enable", "true");
        props.put("mail.smtp.auth.mechanisms", "XOAUTH2");
        props.put("mail.smtp.sasl.enable", "true");
        props.put("mail.smtp.sasl.mechanisms", "XOAUTH2");
        return mailSender;
    }
}
\end{lstlisting}

\subsubsection{\texttt{EmailService}}

Genera un código de 6 dígitos, lo guarda en Redis con una expiración de 10 minutos y lo envía por correo. Al verificar el código, guarda un flag temporal en Redis para permitir el registro.

\begin{lstlisting}[language=Java]
public void sendVerificationCode(String email) {
    String code = String.format("%06d", random.nextInt(1_000_000));
    redis.opsForValue().set(CODE_PREFIX + email, code,
            Duration.ofMinutes(codeTtlMinutes));
    // ... (creacion y envio del correo)
}

public boolean verifyCode(String email, String code) {
    String stored = redis.opsForValue().get(CODE_PREFIX + email);
    if (stored == null || !stored.equals(code)) return false;

    redis.delete(CODE_PREFIX + email);
    redis.opsForValue().set(VERIFIED_PREFIX + email, "true",
            Duration.ofMinutes(verifiedTtlMinutes));
    return true;
}
\end{lstlisting}

\subsubsection{Endpoint de desarrollo \texttt{register-dev}}

Para no perder agilidad en las pruebas, se creó un controlador que \textbf{solo existe en el perfil \texttt{dev}}. Permite crear usuarios directamente sin verificación.

\begin{lstlisting}[language=Java]
@RestController
@RequestMapping("/api/auth")
@Profile("dev")
public class DevAuthController {
    // ... (inyeccion de repositorios)

    @PostMapping("/register-dev")
    public ResponseEntity<AuthResponse> registerDev(@Valid @RequestBody RegisterRequest req) {
        // ... validaciones
        User user = User.builder()
                .email(email)
                .passwordHash(passwordEncoder.encode(req.password()))
                // ...
                .emailVerified(true)
                .build();
        user = userRepository.save(user);
        // ... retorno de tokens
    }
}
\end{lstlisting}

\subsection{Cambios a archivos previos}

\begin{itemize}[leftmargin=*]
  \item \texttt{pom.xml}: se añadieron las dependencias de correo y de Google API Client.
  \item \texttt{application.yml}: se agregó la configuración de \texttt{spring.mail} y el bloque \texttt{notiuam.email}.
  \item \texttt{User.java}: se añadió el campo \texttt{emailVerified}.
  \item \texttt{AuthService.java}: se inyectó \texttt{EmailService} y se añadió la validación en el método \texttt{register}.
\end{itemize}

\subsection{Problemas encontrados y solución}

\subsubsection{Contraseña de aplicación bloqueada por Google Workspace}

La interfaz de contraseñas de aplicación estaba bloqueada por el administrador. \textbf{Solución:} se implementó OAuth2 con un Refresh Token, que no requiere habilitación por parte del administrador.

\subsubsection{Error de credenciales inválidas en el Playground}

Se obtuvo el error \texttt{401: invalid\_client}. \textbf{Solución:} el cliente OAuth debía ser de tipo \textit{Aplicación web}, no de escritorio, para poder usar el Playground.

\subsubsection{\texttt{redirect\_uri\_mismatch}}

Se obtuvo el error \texttt{400: redirect\_uri\_mismatch}. \textbf{Solución:} se agregó la URI \texttt{https://developers.google.com/oauthplayground} a la lista de redirecciones autorizadas del cliente web.

\section{Verificación}

\subsection{Prueba de envío real}

\begin{lstlisting}[language=bash]
curl -s -X POST http://localhost:8081/api/auth/send-verification \
  -H "Content-Type: application/json" \
  -d '{"email":"aaron.r.ramos@cua.uam.mx"}' | jq
\end{lstlisting}

\begin{exito}[Resultado esperado]
\begin{itemize}[leftmargin=*]
  \item Respuesta: \texttt{\{"message": "Codigo enviado"\}}.
  \item Recepción del correo real en la bandeja del usuario.
  \item Log en el backend: \texttt{Correo de verificacion enviado...}.
\end{itemize}
\end{exito}

\subsection{Prueba del atajo de desarrollo}

\begin{lstlisting}[language=bash]
curl -s -X POST http://localhost:8081/api/auth/register-dev \
  -H "Content-Type: application/json" \
  -d '{"email":"test10@cua.uam.mx", "password":"password123", "displayName":"Test", "role":"PROFESSOR"}' | jq
\end{lstlisting}

\begin{exito}[Resultado esperado]
Respuesta \texttt{201 Created} con los tokens de acceso, sin enviar correo.
\end{exito}

\section{Resultado}

\begin{exito}[Estado al cierre]
\begin{itemize}[leftmargin=*]
  \item Envío de correos de verificación reales mediante OAuth2 (XOAUTH2).
  \item Validación de correos antes del registro.
  \item Endpoint de atajo \texttt{register-dev} para pruebas (solo en perfil dev).
\end{itemize}
\end{exito}